\documentclass[12pt,a4paper]{article}

\usepackage[czech]{babel}
\usepackage[T1]{fontenc}
\usepackage{tikz}
\usepackage{amsmath}

\begin{document}

\section*{TicTacToe}

Hra TicTacToe je jednoduchá hra, kterou pravděpodobně zná téměř každý.
Hraje se ve čtvercové síti tvořené třemi řádky a třemi sloupci
(viz obr.~\ref{fig:TicTacToeHraciPole}).

Hráči se pravidelně střídají v umisťování symbolů
$\mathrm{X}$ a $\mathrm{O}$ do jednotlivých polí.
Cílem hráče je jako první vytvořit trojici svých symbolů
v jednom řádku, sloupci nebo na jedné z diagonál.

Tato hra má konečný počet možných průběhů a zakončení.
V tomto textu se pokusíme zjistit, kolik jich skutečně je.

\begin{figure}[h]
    \centering
    \begin{tikzpicture}[scale = 1.5]
        \draw[thick] (1,0) -- (1,3);
        \draw[thick] (2,0) -- (2,3);
        \draw[thick] (0,1) -- (3,1);
        \draw[thick] (0,2) -- (3,2);
    \end{tikzpicture}

    \caption{Hrací pole hry TicTacToe}
    \label{fig:TicTacToeHraciPole}
\end{figure}

Pro usnadnění další práce a lepší orientaci v hrací síti zavedeme
označení jednotlivých polí podobné označení prvků matice.
Pozici budeme zapisovat ve tvaru

\begin{equation}
p(a;b),
\label{eq:OznaceniPozice}
\end{equation}

kde $a$ označuje číslo řádku a $b$ číslo sloupce.
Řádky i sloupce číslujeme od $1$ do $3$, přičemž řádky počítáme
shora dolů a sloupce zleva doprava.

Například $p(2;3)$ označuje pole ve druhém řádku a třetím sloupci
(viz obr.~\ref{fig:TicTacToeVyznaceniP23}).

\begin{figure}[h]
    \begin{center}
        \begin{tikzpicture}[scale = 1.5]
            \draw[thick] (1,0) -- (1,3);
            \draw[thick] (2,0) -- (2,3);
            \draw[thick] (0,1) -- (3,1);
            \draw[thick] (0,2) -- (3,2);

            \fill (2.5,1.5) circle (2pt);
        \end{tikzpicture}

        \caption{Vyznačení pozice $p(2;3)$}
        \label{fig:TicTacToeVyznaceniP23}
    \end{center}
\end{figure}

Dále si označíme hráče.
Hráč hrající $\mathrm{X}$ bude začínat a označovat jej budeme $\mathrm{X}$.
Hráč hrající $\mathrm{O}$ bude hrát vždy jako druhý a označovat jej budeme $\mathrm{O}$.

Nejprve musíme zjistit, kolik různých posloupností tahů lze v TicTacToe vytvořit.
Prozatím budeme předpokládat, že se hra vždy odehraje až do obsazení všech devíti polí.

Hráč $\mathrm{X}$ může svůj symbol vložit na libovolnou z $9$ pozic.
Poté již zbývá $8$ pozic pro hráče $\mathrm{O}$.
Po jeho tahu zbývá pouze $7$ volných pozic a každým dalším tahem se počet
volných míst na hrací ploše sníží o jedna.

Pro první tah hráče $\mathrm{X}$ máme tedy $9$ možností umístění.
Pro každou z těchto možností existuje $8$ možných tahů hráče $\mathrm{O}$,
pro každý z nich dalších $7$ možností hráče $\mathrm{X}$ a tak dále.

Celkový počet možných posloupností devíti tahů můžeme zapsat jako

\begin{equation}
P(9) = 9! = 9\cdot8\cdot7\cdot6\cdot5\cdot4\cdot3\cdot2\cdot1 = 362\,880.
\label{eq:PocetVsechPoradi}
\end{equation}

Tento počet však stále obsahuje hry, které se od sebe liší pouze otočením
nebo zrcadlením hrací plochy.

Například následující dvě posloupnosti tahů se od sebe liší pouze otočením
hrací plochy o $90^\circ$.

\begin{figure}[h]
    \begin{center}
        \begin{tikzpicture}[scale = 1.2]

            \begin{scope}
                \draw[thick] (1,0) -- (1,3);
                \draw[thick] (2,0) -- (2,3);
                \draw[thick] (0,1) -- (3,1);
                \draw[thick] (0,2) -- (3,2);

                \node at (0.5,2.5) {\Large $X_1$};
                \node at (1.5,1.5) {\Large $O_1$};
                \node at (1.5,0.5) {\Large $O_2$};
                \node at (2.5,0.5) {\Large $X_2$};
            \end{scope}

            \begin{scope}[xshift = 5cm]
                \draw[thick] (1,0) -- (1,3);
                \draw[thick] (2,0) -- (2,3);
                \draw[thick] (0,1) -- (3,1);
                \draw[thick] (0,2) -- (3,2);

                \node at (2.5,2.5) {\Large $X_1$};
                \node at (1.5,1.5) {\Large $O_1$};
                \node at (0.5,1.5) {\Large $O_2$};
                \node at (0.5,0.5) {\Large $X_2$};
            \end{scope}

            \draw[->,thick] (3.4,1.5) -- (4.6,1.5)
                node[midway,above] {$90^\circ$};

        \end{tikzpicture}

        \caption{Dvě posloupnosti tahů lišící se pouze otočením hrací plochy}
        \label{fig:TicTacToeRotaceHraciPlochy}
    \end{center}
\end{figure}

Hrací plocha má celkem 8 symetrií. Čtyři z nich jsou osové souměrnosti podle:
\begin{itemize}
    \item podle vodorovné osy procházející středem hrací plochy,
    \item podle svislé osy procházející středem hrací plochy,
    \item diagonály vedoucí poli $p(1;1)$ a $p(3;3)$,
    \item diagonály vedoucí poli $p(1;3)$ a $p(3;1)$.
\end{itemize}

Tyto osové souměrnosti jsou znázorněny na obr.~\ref{fig:TicTacToeOsovaSymetrieHry}.

\begin{figure}[h]
    \begin{center}
        \begin{tikzpicture}[scale = 1.5]
            \draw[thick] (1,0) -- (1,3);
            \draw[thick] (2,0) -- (2,3);
            \draw[thick] (0,1) -- (3,1);
            \draw[thick] (0,2) -- (3,2);

            \draw[thick,red] (1.5,2.7) -- (1.5,0.3) node[below right] {$o_1$};
            \draw[thick,red] (0.3,1.5) -- (2.7,1.5) node[below right] {$o_2$};
            \draw[thick,red] (0.3,2.7) -- (2.7,0.3) node[below right] {$o_3$};
            \draw[thick,red] (0.3,0.3) -- (2.7,2.7) node[below right] {$o_4$};
        \end{tikzpicture}
    \end{center}

    \caption{Osové souměrnosti hrací plochy}
    \label{fig:TicTacToeOsovaSymetrieHry}
\end{figure}

Zbývající čtyři symetrie jsou otočení hrací plochy kolem jejího středu o:
\begin{itemize}
    \item $90^\circ$,
    \item $180^\circ$,
    \item $270^\circ$,
    \item $360^\circ = 0^\circ$
\end{itemize}

Otočení o $0^\circ$ ponechává hrací plochu beze změny a představuje tedy
původní polohu hrací plochy. Ostatní tři otočení mění pozice jednotlivých polí,
ale zachovávají jejich vzájemné uspořádání.

Máme tedy celkem $8$ různých symetrií hrací plochy:
$4$ osové souměrnosti a $4$ otočení.

Nyní můžeme tyto symetrie využít při počítání různých posloupností tahů.
Pokud lze jednu posloupnost tahů získat z jiné pouze otočením nebo osovou
souměrností celé hrací plochy, budeme je považovat za stejnou hru.

Každá úplná posloupnost devíti tahů má právě $8$ takových symetrických podob.
Žádná z nenulových symetrií totiž nemůže úplnou posloupnost tahů převést samu
na sebe. V takovém případě by musela ponechat na stejném místě první tah,
druhý tah a postupně všechny další tahy. Protože je po devíti tazích obsazeno
všech $9$ polí, musela by tato symetrie ponechat na místě všechna pole hrací
plochy, což splňuje pouze otočení o $0^\circ$.

Počet různých posloupností tahů po odstranění symetrických variant je tedy

\begin{equation}
    \frac{9!}{8} = \frac{362\,880}{8} = 45\,360.
\end{equation}

Máme tedy $45\,360$ různých posloupností devíti tahů, pokud posloupnosti,
které se od sebe liší pouze otočením nebo osovou souměrností hrací plochy,
považujeme za stejné.

Tento výsledek však stále není konečný.
Doposud jsme totiž předpokládali, že každá hra obsahuje právě $9$ tahů.
Ve skutečnosti může hra skončit dříve, pokud některý z hráčů vytvoří
vítěznou trojici.

Protože první vítězná trojice může vzniknout už v pátém tahu, 
musíme spočítat každou z variant ukončení v tazích od 5. do 9. zvlášť 
a poté tyto varianty sečíst.

Nejprve začneme určovat počet her končících 5. tahem.
Po pěti tazích je vítězem $\mathrm{X}$ a tedy vytvořil vítěznou trojici, 
zatímco $\mathrm{O}$ má položené pouze dva tokeny.

Vítěznými trojicemi jsou buď tři řádky, tři sloupce, nebo dvě diagonály. 
Vítězných trojic je tedy:
\[
    3+3+2 = 8
\]
V těchto jednotlivých výherních pozicích může $\mathrm{X}$ položit tokeny:
\[
    3! = 3\cdot2\cdot1 = 6
\]
způsoby.

Po obsazení tří vítězných míst zbývá $\mathrm{O}$ 6 polí, kde může hrát. 
A protože má dva tokeny, tak první může umístit na 6 volných míst 
a druhý token už jen na 5. Tedy:
\[
    6\cdot5 = 30
\]
způsobů pořadí umístění tokenů.

Dosadíme a dostáváme
\[
    8\cdot3!\cdot6\cdot5 = 8\cdot6\cdot30 = 1\,440.
\]
Her končících po pátém tahu je tedy celkem $1\,440$, 
tento počet ale opět obsahuje i symetrické varianty. 
Musíme tedy výsledek o tyto symetrie ošetřit.
\begin{equation}
    \frac{1\,440}{8} = 180
\end{equation}
Celkem tedy můžeme mít $180$ různých her, 
které byly ukončeny vítězstvím $\mathrm{X}$ po pátém tahu.

Nyní určíme počet her končících 6. tahem.
Po šesti tazích je vítězem $\mathrm{O}$, který vytvořil vítěznou trojici, 
zatímco $\mathrm{X}$ má položené tři tokeny.

Stejně jako v předchozím případě může vítězná trojice hráče $\mathrm{O}$ vzniknout 8 způsoby 
a své tři tokeny může v této trojici položit:
\[
    3! = 6
\]
způsoby.

Hráč $\mathrm{X}$ může své tři tokeny položit na tři ze zbývajících 6 polí. 
První token může položit na 6 polí, druhý na 5 a třetí na 4. Tedy:
\[
    6\cdot5\cdot4 = 120
\]
způsoby.

Dostáváme tedy:
\[
    8\cdot3!\cdot6\cdot5\cdot4 = 5\,760
\]
možných her.

V tomto počtu jsou však zahrnuty také hry, ve kterých již $\mathrm{X}$ zvítězil v 5. tahu. 
Tyto hry tedy musíme odečíst.

Aby mohl $\mathrm{X}$ zvítězit v 5. tahu a zároveň měl $\mathrm{O}$ po 6. tahu vítěznou trojici,
musí být jejich vítězné trojice navzájem oddělené. 
To je možné pouze tehdy, pokud jsou obě vítězné trojice tvořeny dvěma různými řádky nebo dvěma různými sloupci.

Pro vítěznou trojici hráče $\mathrm{O}$ máme tedy 6 možností, 
protože musí jít o jeden ze tří řádků nebo jeden ze tří sloupců. 
Pro každou z těchto možností existují 2 vítězné trojice hráče $\mathrm{X}$, 
které s ní nemají žádné společné pole. 
Tokeny obou hráčů mohou být ve svých trojicích položeny $3!$ způsoby. 
Dostáváme tedy:
\[
    6\cdot2\cdot3!\cdot3! = 432
\]
her, které již skončily vítězstvím hráče $\mathrm{X}$ v 5. tahu.

Počet her skutečně končících 6. tahem je tedy:
\[
    5\,760-432 = 5\,328
\]

Po ošetření symetrií dostáváme:
\begin{equation}
    \frac{5\,328}{8} = 666
\end{equation}

Celkem tedy můžeme mít $666$ různých her, 
které byly ukončeny vítězstvím $\mathrm{O}$ po šestém tahu.

Nyní určíme počet her končících 7. tahem.
Po sedmi tazích je vítězem opět $\mathrm{X}$. 
Hráč $\mathrm{X}$ má v tomto okamžiku položené čtyři tokeny, 
přičemž tři z nich musí tvořit vítěznou trojici a jeden token leží mimo ni.

Vítěznou trojici můžeme opět vybrat 8 způsoby a 
čtvrtý token hráče $\mathrm{X}$ může ležet na jednom ze zbývajících 6 polí.

Aby hra skončila skutečně až 7. tahem, musí být posledním položeným tokenem hráče 
$\mathrm{X}$ jeden ze tří tokenů tvořících vítěznou trojici. 
Pro poslední token tedy máme 3 možnosti 
a zbývající tři tokeny hráče $\mathrm{X}$ můžeme položit:
\[
    3! = 6
\]
způsoby.

Pro každé rozmístění tokenů hráče $\mathrm{X}$ tedy máme:
\[
    3\cdot3! = 18
\]
možných pořadí jeho tahů.

Nyní musíme určit možnosti hráče $\mathrm{O}$. 
Pokud je vítěznou trojicí hráče $\mathrm{X}$ řádek nebo sloupec, 
máme 6 možností vítězné trojice a 6 možností umístění čtvrtého tokenu. 
Celkem tedy:
\[
    6\cdot6 = 36
\]
možností.

Ve zbývajících pěti polích se v tomto případě nachází jedna celá vítězná trojice, 
kterou by mohl vytvořit hráč $\mathrm{O}$. 
Hráč $\mathrm{O}$ může své tři tokeny obecně položit:
\[
    5\cdot4\cdot3 = 60
\]
způsoby. Z toho je ale $3! = 6$ možností, 
při kterých by již vytvořil vítěznou trojici v 6. tahu. 
Zbývá tedy:
\[
    60-6 = 54
\]
možností.

Pokud je vítěznou trojicí hráče $\mathrm{X}$ diagonála, 
máme 2 možnosti vítězné trojice a opět 6 možností umístění čtvrtého tokenu. 
Tedy:
\[
    2\cdot6 = 12
\]
možností.

V tomto případě se mezi zbývajícími pěti poli nenachází žádná celá vítězná trojice. 
Hráč $\mathrm{O}$ tedy může své tokeny položit všemi:
\[
    5\cdot4\cdot3 = 60
\]
způsoby.

Celkový počet her končících 7. tahem je tedy:
\[
    36\cdot18\cdot54+12\cdot18\cdot60 = 47\,952
\]

Po ošetření symetrií dostáváme:
\begin{equation}
    \frac{47\,952}{8} = 5\,994
\end{equation}

Celkem tedy můžeme mít $5\,994$ různých her, 
které byly ukončeny vítězstvím $\mathrm{X}$ po sedmém tahu.

Nyní určíme počet her končících 8. tahem.
Po osmi tazích je vítězem $\mathrm{O}$. 
Hráč $\mathrm{O}$ má položené čtyři tokeny, 
přičemž tři z nich tvoří vítěznou trojici a jeden token leží mimo ni.

Stejně jako v předchozím případě může být posledním tokenem hráče 
$\mathrm{O}$ jeden ze tří tokenů tvořících vítěznou trojici. 
Zbývající tři tokeny může položit $3!$ způsoby. 
Pro každé rozmístění jeho čtyř tokenů máme tedy:
\[
    3\cdot3! = 18
\]
možných pořadí tahů.

Opět musíme rozlišit, zda je vítěznou trojicí řádek nebo sloupec, nebo zda jde o diagonálu.

Pokud je vítěznou trojicí hráče $\mathrm{O}$ řádek nebo sloupec, máme:
\[
    6\cdot6 = 36
\]
možností rozmístění jeho čtyř tokenů.

Ve zbývajících pěti polích se nachází jedna vítězná trojice, 
kterou by mohl vytvořit hráč $\mathrm{X}$. 
Protože hráč $\mathrm{X}$ obsadí čtyři z těchto pěti polí, 
musí jediné neobsazené pole ležet právě v této vítězné trojici. 
Pro neobsazené pole máme tedy 3 možnosti.

Čtyři tokeny hráče $\mathrm{X}$ mohou být následně položeny:
\[
    4! = 24
\]
způsoby.

Pokud je vítěznou trojicí hráče $\mathrm{O}$ diagonála, máme:
\[
    2\cdot6 = 12
\]
možností rozmístění jeho čtyř tokenů.

Mezi zbývajícími pěti poli v tomto případě není žádná vítězná trojice. 
Jediné neobsazené pole tedy může být kterékoliv z těchto 5 polí.

Celkový počet her končících 8. tahem je tedy:
\[
    \left(36\cdot3+12\cdot5\right)\cdot18\cdot4! = 72\,576
\]

Po ošetření symetrií dostáváme:
\begin{equation}
    \frac{72\,576}{8} = 9\,072
\end{equation}

Celkem tedy můžeme mít $9\,072$ různých her, 
které byly ukončeny vítězstvím $\mathrm{O}$ po osmém tahu.

Zbývá určit počet her končících až 9. tahem.
Z původních $9! = 362\,880$ posloupností musíme odečíst všechny posloupnosti, 
které by ve skutečnosti skončily již po 5., 6., 7. nebo 8. tahu.

Hra ukončená po 5. tahu je v původním počtu $9!$ obsažena vícekrát.
Po ukončení hry totiž zbývají 4 neobsazená pole, 
která mohou být v původní devítitahové posloupnosti seřazena:
\[
    4! = 24
\]
způsoby.

Stejně tak hra končící po 6. tahu je započítána $3!$ krát, 
hra končící po 7. tahu $2!$ krát a hra končící po 8. tahu $1!$ krát.

Počet posloupností, které skutečně pokračují až k 9. tahu, je tedy:
\[
    9!-1\,440\cdot4!-5\,328\cdot3!-47\,952\cdot2!-72\,576\cdot1!
\]
tedy:
\[
    362\,880-34\,560-31\,968-95\,904-72\,576 = 127\,872
\]

Tento počet obsahuje jak hry, ve kterých hráč $\mathrm{X}$ vytvoří vítěznou trojici až v 9. tahu,
 tak hry, ve kterých ani jeden z hráčů nevytvoří vítěznou trojici a hra skončí remízou.

Po ošetření symetrií dostáváme:
\begin{equation}
    \frac{127\,872}{8} = 15\,984
\end{equation}

Celkem tedy můžeme mít $15\,984$ různých her, které byly ukončeny po devátém tahu.

Nyní můžeme sečíst všechny možnosti ukončení hry (výsledky rovnic 4--8):
\[
    180+666+5\,994+9\,072+15\,984 = 31\,896
\]

Celkově tedy existuje $31\,896$ různých her TicTacToe, pokud hry, 
které se od sebe liší pouze otočením nebo osovou souměrností hrací plochy, 
považujeme za stejné.

Doposud jsme počítali jednotlivé hry podle pořadí, 
ve kterém byly tokeny na hrací plochu pokládány.
Dvě různé posloupnosti tahů však mohou vést ke stejné pozici na hrací ploše.

Nyní se tedy můžeme podívat na jinou otázku. 
Kolik různých pozic může během hry TicTacToe vzniknout?

Na začátku hry máme pouze jednu možnou pozici, a to prázdnou hrací plochu.
Po prvním tahu má hráč $\mathrm{X}$ položený jeden token. 
Ten může umístit na libovolné z 9 polí. Počet pozic po prvním tahu je tedy:
\[
    \binom{9}{1}=9
\]

Po druhém tahu má hráč $\mathrm{X}$ položený jeden token a hráč $\mathrm{O}$ také jeden token. 
Hráč $\mathrm{X}$ může obsadit jedno z 9 polí a hráč $\mathrm{O}$ následně jedno ze zbývajících 8 polí. 
Dostáváme tedy:
\[
    9\cdot8=72
\]
různých pozic.

Po třetím tahu má hráč $\mathrm{X}$ položené dva tokeny a hráč $\mathrm{O}$ jeden token.
Na rozdíl od počítání jednotlivých her nás nyní nezajímá pořadí, 
ve kterém hráč $\mathrm{X}$ své dva tokeny položil. 
Z 9 polí tedy pouze vybíráme dvě pole, která budou obsazena tokeny hráče $\mathrm{X}$.

Počet možností pro hráče $\mathrm{X}$ je:
\[
    \binom{9}{2}=36
\]

Hráči $\mathrm{O}$ poté zbývá 7 možných polí. Celkem tedy dostáváme:
\[
    \binom{9}{2}\cdot7=36\cdot7=252
\]
různých pozic.

Po čtvrtém tahu mají oba hráči položené dva tokeny. 
Nejprve vybereme dvě z 9 polí pro hráče $\mathrm{X}$ 
a následně dvě ze zbývajících 7 polí pro hráče $\mathrm{O}$. 
Tedy:
\[
    \binom{9}{2}\cdot\binom{7}{2}=36\cdot21=756
\]
různých pozic.

Po pátém tahu má hráč $\mathrm{X}$ položené tři tokeny a hráč $\mathrm{O}$ dva tokeny. 
Pro hráče $\mathrm{X}$ vybereme tři z 9 polí a pro hráče $\mathrm{O}$ dvě ze zbývajících 6 polí:
\[
    \binom{9}{3}\cdot\binom{6}{2}=84\cdot15=1\,260
\]

Po pátém tahu tedy existuje $1\,260$ různých pozic. 
V tomto případě zatím nemusíme žádné pozice odečítat, 
protože pátý tah je prvním tahem, ve kterém může některý z hráčů vyhrát. 
Pozice, ve které hráč $\mathrm{X}$ vytvořil vítěznou trojici právě v tomto tahu, 
je tedy stále platnou konečnou pozicí hry.

Po šestém tahu mají oba hráči položené tři tokeny. 
Pokud bychom zatím nebrali v úvahu možné dřívější ukončení hry, 
dostali bychom:
\[
    \binom{9}{3}\cdot\binom{6}{3}=84\cdot20=1\,680
\]
různých pozic.

Ne všechny tyto pozice však mohou ve skutečné hře nastat. 
Pokud tři tokeny hráče $\mathrm{X}$ tvoří vítěznou trojici, 
hra skončila již po pátém tahu a hráč $\mathrm{O}$ již svůj třetí token položit nemohl.

Hráč $\mathrm{X}$ může vytvořit jednu z 8 vítězných trojic. 
Hráč $\mathrm{O}$ potom může své tři tokeny rozmístit na tři ze zbývajících 6 polí:
\[
    8\cdot\binom{6}{3}=8\cdot20=160
\]
způsoby.

Tyto pozice musíme od původního počtu odečíst:
\[
    1\,680-160=1\,520
\]

Po šestém tahu tedy může skutečně nastat $1\,520$ různých pozic.

Po sedmém tahu má hráč $\mathrm{X}$ položené čtyři tokeny a hráč $\mathrm{O}$ tři tokeny. 
Bez omezení bychom dostali:
\[
    \binom{9}{4}\cdot\binom{5}{3}=126\cdot10=1\,260
\]
různých pozic.

Nyní však musíme odstranit pozice, 
ve kterých již hráč $\mathrm{O}$ vytvořil vítěznou trojici v šestém tahu. 
Protože má hráč $\mathrm{O}$ právě tři tokeny, 
musí tyto tři tokeny tvořit jednu z 8 vítězných trojic.

Pro každou vítěznou trojici hráče $\mathrm{O}$ může hráč $\mathrm{X}$ obsadit 
čtyři ze zbývajících 6 polí. Počet neplatných pozic je tedy:
\[
    8\cdot\binom{6}{4}=8\cdot15=120
\]

Počet možných pozic po sedmém tahu je tedy:
\[
    1\,260-120=1\,140
\]

Po osmém tahu mají oba hráči položené čtyři tokeny. 
Bez omezení bychom dostali:
\[
    \binom{9}{4}\cdot\binom{5}{4}=126\cdot5=630
\]
různých pozic.

Musíme však odstranit pozice, ve kterých hráč $\mathrm{X}$ vytvořil vítěznou trojici 
již před osmým tahem.

Hráč $\mathrm{X}$ má čtyři tokeny. 
Jeho vítěznou trojici můžeme vybrat 8 způsoby 
a čtvrtý token může ležet na jednom ze zbývajících 6 polí. 
Dostáváme tedy:
\[
    8\cdot6=48
\]
možných rozmístění tokenů hráče $\mathrm{X}$ obsahujících vítěznou trojici.

Pro každé takové rozmístění může hráč $\mathrm{O}$ obsadit čtyři z 5 zbývajících polí:
\[
    \binom{5}{4}=5
\]
způsoby.

Počet neplatných pozic je tedy:
\[
    48\cdot5=240
\]

Počet možných pozic po osmém tahu je:
\[
    630-240=390
\]

Po devátém tahu má hráč $\mathrm{X}$ položeno pět tokenů 
a hráč $\mathrm{O}$ čtyři tokeny. 
Hrací plocha je tedy již zcela zaplněna.

Bez omezení můžeme vybrat pět z 9 polí pro hráče $\mathrm{X}$ 
a zbývající čtyři pole již automaticky patří hráči $\mathrm{O}$:
\[
    \binom{9}{5}=126
\]
možnými způsoby.

Musíme však odstranit pozice, ve kterých hráč $\mathrm{O}$ již vytvořil 
vítěznou trojici v osmém tahu.

Stejně jako v předchozím případě může hráč $\mathrm{O}$ vytvořit jednu z 
8 vítězných trojic a svůj čtvrtý token umístit na jedno ze zbývajících 6 polí. 
Takových pozic je:
\[
    8\cdot6=48
\]

Počet možných konečných pozic po devátém tahu je tedy:
\[
    126-48=78
\]

Nyní můžeme sečíst počty všech pozic, které mohou během hry nastat. 
Započítáme také počáteční prázdnou hrací plochu:
\begin{equation}
    1+9+72+252+756+1\,260+1\,520+1\,140+390+78=5\,478.
\end{equation}

Celkem tedy může ve hře TicTacToe nastat $5\,478$ různých pozic,
 pokud započítáme také počáteční prázdnou hrací plochu. 
 Pokud bychom počítali pouze pozice vzniklé po provedení 
 alespoň jednoho tahu, dostali bychom:
\[
    5\,478-1=5\,477
\]
různých pozic.

Tento počet však opět nerozlišuje pozice podle jejich symetrie. 
Mohlo by se tedy zdát, že stejně jako v případě jednotlivých her stačí výsledný počet vydělit 8.

U pozic to však již není možné.

Například pozice, ve které hráč $\mathrm{X}$ položí svůj první token do středu hrací plochy, 
zůstane při každém otočení i osové souměrnosti stejná. 
Naproti tomu pozice, ve které je token položen v rohu, 
může mít několik různých symetrických podob.

Jednotlivé pozice tedy nemají vždy 8 různých symetrických variant. 
Počet různých pozic po zohlednění symetrií proto nemůžeme získat 
jednoduchým vydělením čísla $5\,478$ osmi.

U jednotlivých pozic tedy musíme postupovat jiným způsobem.

Po prvním tahu jsme získali celkem 9 různých pozic. Pokud ale zohledníme symetrie hrací plochy, existují pouze tři navzájem různé možnosti. Hráč $\mathrm{X}$ může položit svůj token:
\begin{itemize}
    \item do středu hrací plochy,
    \item do jednoho z rohů,
    \item doprostřed jedné ze stran.
\end{itemize}

Všechny čtyři rohy lze pomocí otočení převést jeden na druhý a stejně tak lze mezi sebou převádět všechna čtyři pole ležící uprostřed stran. Středové pole zůstává při všech symetriích na stejném místě.

Z původních 9 pozic tedy po prvním tahu dostáváme pouze:
\[
    3
\]
různé pozice po zohlednění symetrií.

Po druhém tahu máme celkem 72 pozic. 
Nyní musíme pro každou ze tří možností prvního tahu určit, 
kolik různých možností má hráč $\mathrm{O}$.

Pokud hráč $\mathrm{X}$ zahrál do středu, 
jsou všechna čtyři rohová pole navzájem symetrická 
a stejně tak všechna čtyři pole uprostřed stran. 
Hráč $\mathrm{O}$ má tedy pouze 2 navzájem různé možnosti tahu.

Pokud hráč $\mathrm{X}$ zahrál do rohu, 
můžeme zbývající pole rozdělit podle jejich polohy vůči tomuto rohu. 
Po zohlednění symetrií, které zachovávají původní roh na jeho místě, 
dostáváme 5 různých možností tahu hráče $\mathrm{O}$.

Stejně tak pokud hráč $\mathrm{X}$ zahrál doprostřed strany, 
dostáváme pro hráče $\mathrm{O}$ opět 5 navzájem různých možností.

Celkem tedy po druhém tahu existuje:
\[
    2+5+5=12
\]
různých pozic po zohlednění symetrií.

Od třetího tahu již začíná být podobné ruční rozdělování pozic 
na jednotlivé případy značně nepřehledné. 
Proto si zavedeme jednotný způsob, kterým budeme symetrické pozice porovnávat.

Každou pozici můžeme pomocí osmi symetrií hrací plochy převést až na 8 různých podob. 
Z těchto podob vždy vybereme jednu, kterou budeme považovat za reprezentanta dané pozice.

Pro jednoznačný výběr reprezentanta si můžeme jednotlivá pole hrací plochy 
zapsat postupně po řádcích a každému poli přiřadit hodnotu:
\[
    0=\text{prázdné pole}, \qquad 1=\mathrm{O}, \qquad 2=\mathrm{X}.
\]

Každou pozici tak můžeme zapsat jako posloupnost devíti čísel. Například pozici:
\[
\begin{array}{c|c|c}
    \mathrm{X} & & \\
    \hline
    & \mathrm{O} & \\
    \hline
    & & \mathrm{X}
\end{array}
\]
můžeme zapsat jako:
\[
    200010002.
\]

Pro každou pozici vytvoříme všech 8 podob získaných pomocí symetrií hrací plochy 
a jako jejího reprezentanta vybereme vždy číselně nejmenší z těchto zápisů.

Pokud mají dvě pozice stejného reprezentanta, znamená to, 
že lze jednu z nich převést na druhou pomocí některé ze symetrií hrací plochy. 
Považujeme je tedy za stejnou pozici.

Tímto způsobem můžeme postupně projít všechny možné pozice, 
které jsme získali v předchozí části, a započítat každého reprezentanta pouze jednou.

Po prvním tahu jsme již dostali:
\[
    3
\]
různé pozice a po druhém tahu:
\[
    12.
\]

Stejným postupem pro další tahy dostáváme:
\[
\begin{array}{c|c|c}
    \text{Počet tahů} & \text{Pozic bez zohlednění symetrií} & \text{Pozic po zohlednění symetrií} \\
    \hline
    0 & 1 & 1 \\
    1 & 9 & 3 \\
    2 & 72 & 12 \\
    3 & 252 & 38 \\
    4 & 756 & 108 \\
    5 & 1\,260 & 174 \\
    6 & 1\,520 & 204 \\
    7 & 1\,140 & 153 \\
    8 & 390 & 57 \\
    9 & 78 & 15
\end{array}
\]

Nyní můžeme všechny navzájem různé pozice sečíst:
\begin{equation}
    1+3+12+38+108+174+204+153+57+15=765.
\end{equation}

Celkem tedy může během hry TicTacToe vzniknout $765$ navzájem různých pozic, 
pokud pozice lišící se pouze otočením nebo osovou souměrností hrací plochy považujeme za stejné.

Do tohoto počtu je opět započítána i počáteční prázdná hrací plocha. 
Pokud bychom ji nezapočítali, dostáváme:
\[
    765-1=764
\]
různých pozic vzniklých po provedení alespoň jednoho tahu.

Tím jsme určili jak počet různých průběhů hry, tak počet různých pozic,
které mohou během hry vzniknout. 
Zůstává však ještě jedna přirozená otázka. 
Pokud oba hráči hrají optimálně, může některý z nich vždy zvítězit, 
nebo může druhý hráč vítězství zabránit?

Abychom mohli tuto otázku zodpovědět, můžeme se na hru podívat jako na strom možností.

Na začátku máme prázdnou hrací plochu. Z ní vede 9 možných prvních tahů hráče $\mathrm{X}$.
Z každé takto vzniklé pozice vede následně až 8 možných tahů hráče $\mathrm{O}$ a tímto způsobem můžeme pokračovat až do okamžiku, kdy některý z hráčů zvítězí nebo hra skončí remízou.

Takovou strukturu nazýváme herní strom.

Každý vrchol tohoto stromu představuje jednu pozici hry 
a hrany mezi jednotlivými vrcholy představují možné tahy hráčů.

Na začátku však není nutné rozlišovat všech 9 možností prvního tahu.
Jak jsme již ukázali při práci se symetriemi, 
po prvním tahu existují pouze 3 navzájem různé pozice.

Hráč $\mathrm{X}$ může začít:
\begin{itemize}
    \item ve středu hrací plochy,
    \item v jednom z rohů,
    \item uprostřed jedné ze stran.
\end{itemize}

Ostatní možnosti lze na některou z těchto tří převést pomocí otočení 
nebo osové souměrnosti hrací plochy.

Nyní potřebujeme určit, která z těchto možností vede při optimální hře k vítězství, 
prohře nebo remíze.

Jednotlivým konečným pozicím proto přiřadíme hodnotu podle výsledku hry 
z pohledu hráče $\mathrm{X}$:
\[
    v=
    \begin{cases}
        1, & \text{pokud zvítězí } \mathrm{X},\\
        0, & \text{pokud hra skončí remízou},\\
        -1, & \text{pokud zvítězí } \mathrm{O}.
    \end{cases}
\]

Pokud je na tahu hráč $\mathrm{X}$, bude se snažit vybrat takový tah, 
který vede k co nejvyšší možné hodnotě. Hodnotu dané pozice tedy určíme jako:
\[
    v(P)=\max\{v(P_1),v(P_2),\ldots,v(P_n)\},
\]
kde $P_1,P_2,\ldots,P_n$ jsou všechny pozice, 
které mohou vzniknout jedním tahem hráče $\mathrm{X}$.

Naopak hráč $\mathrm{O}$ se snaží hodnotu z pohledu hráče $\mathrm{X}$ co nejvíce snížit. 
Pokud je tedy na tahu hráč $\mathrm{O}$, dostáváme:
\[
    v(P)=\min\{v(P_1),v(P_2),\ldots,v(P_n)\}.
\]

Tímto způsobem můžeme začít u konečných pozic hry a postupně určovat hodnoty předchozích pozic,
dokud nedojdeme až k počáteční prázdné hrací ploše.

Tento postup nazýváme algoritmus minimax.

Pokud má například hráč $\mathrm{X}$ možnost přejít do pozic s hodnotami:
\[
    -1,\qquad 0,\qquad 1,
\]
vybere si možnost s nejvyšší hodnotou:
\[
    \max\{-1,0,1\}=1.
\]

Naopak pokud je ve stejné situaci na tahu hráč $\mathrm{O}$, vybere možnost:
\[
    \min\{-1,0,1\}=-1.
\]

Hodnotu pozice tedy neurčuje pouze to, zda z ní existuje cesta k vítězství.
Důležité je také to, zda může soupeř této cestě zabránit.

Pokud tímto způsobem vyhodnotíme celý herní strom TicTacToe, 
dostáváme pro počáteční prázdnou hrací plochu hodnotu:
\[
    v(P_0)=0.
\]

To znamená, že pokud oba hráči hrají optimálně, hra vždy skončí remízou.

Hráč $\mathrm{X}$ si tedy nemůže vynutit vítězství bez ohledu na tahy soupeře.
Stejně tak si však hráč $\mathrm{O}$ nemůže vynutit vítězství 
proti optimálně hrajícímu hráči $\mathrm{X}$.

Výsledek hry při optimální strategii obou hráčů je tedy:
\[
    \mathrm{X}\;-\;\mathrm{O}=\text{remíza}.
\]

Nyní se můžeme podívat na jednotlivé možnosti prvního tahu hráče $\mathrm{X}$ podrobněji.

Jak jsme již ukázali, po zohlednění symetrií existují pouze tři různé možnosti prvního tahu:
\begin{itemize}
    \item střed hrací plochy,
    \item roh hrací plochy,
    \item střed strany hrací plochy.
\end{itemize}

Pokud oba hráči hrají optimálně, mají všechny tři možnosti stejnou hodnotu:
\[
    v(P)=0.
\]

Bez ohledu na první tah hráče $\mathrm{X}$ tedy může hráč 
$\mathrm{O}$ správnou hrou zabránit jeho vítězství a dovést hru k remíze.

Jednotlivé první tahy se však liší v tom, 
kolik správných odpovědí má hráč $\mathrm{O}$ k dispozici.

Nejprve předpokládejme, že hráč $\mathrm{X}$ zahraje do středu hrací plochy.

Zbývajících 8 polí můžeme vzhledem k symetrii rozdělit na dva typy:
\begin{itemize}
    \item 4 rohová pole,
    \item 4 pole ležící uprostřed stran.
\end{itemize}

Pokud hráč $\mathrm{O}$ odpoví tahem do některého z rohů, 
může při dalším optimálním hraní vynutit remízu. 
Pokud však zahraje doprostřed některé ze stran, 
může hráč $\mathrm{X}$ při správném pokračování vynutit vítězství.

Hráč $\mathrm{O}$ má tedy po tahu hráče $\mathrm{X}$ do středu:
\[
    4
\]
tahy vedoucí při optimální další hře k remíze a:
\[
    4
\]
tahy, po kterých může hráč $\mathrm{X}$ vynutit vítězství.

Nyní předpokládejme, že hráč $\mathrm{X}$ zahraje do některého z rohů.

V tomto případě existuje pro hráče $\mathrm{O}$ pouze jediný tah, 
kterým může při optimální další hře zabránit vynucenému vítězství hráče $\mathrm{X}$. 
Musí zahrát do středu hrací plochy.

Z 8 možných odpovědí hráče $\mathrm{O}$ tedy pouze:
\[
    1
\]
vede při správném pokračování k remíze, zatímco po zbývajících:
\[
    7
\]
tazích může hráč $\mathrm{X}$ vynutit vítězství.

Nakonec předpokládejme, že hráč $\mathrm{X}$ zahraje doprostřed některé ze stran.

Také v tomto případě má hráč $\mathrm{O}$ několik možností, 
kterými může zabránit vynucenému vítězství hráče $\mathrm{X}$.

Pokud například hráč $\mathrm{X}$ zahraje na pole $p(1;2)$,
může hráč $\mathrm{O}$ odpovědět na:
\[
    p(1;1),\qquad p(1;3),\qquad p(2;2),\qquad p(3;2).
\]
Tyto čtyři tahy vedou při optimální další hře k remíze.

Po zbývajících čtyřech možnostech:
\[
    p(2;1),\qquad p(2;3),\qquad p(3;1),\qquad p(3;3)
\]
může hráč $\mathrm{X}$ při správném pokračování vynutit vítězství.

Pro jednotlivé typy prvního tahu tedy dostáváme:
\[
\begin{array}{c|c|c}
    \text{První tah } \mathrm{X} & \text{Odpovědi vedoucí k remíze} & \text{Odpovědi umožňující výhru } \mathrm{X} \\
    \hline
    \text{střed} & 4 & 4 \\
    \text{roh} & 1 & 7 \\
    \text{střed strany} & 4 & 4
\end{array}
\]
Z pohledu optimální hry jsou tedy všechny tři první tahy rovnocenné, 
protože proti správně hrajícímu soupeři vedou k remíze.

Pokud však předpokládáme, že soupeř nemusí vždy zahrát optimálně,
je nejzajímavějším prvním tahem roh hrací plochy. 
Hráč $\mathrm{O}$ má totiž v takovém případě pouze jedinou z 8 možných odpovědí, 
která mu umožňuje zabránit vynucenému vítězství hráče $\mathrm{X}$.

Pravděpodobnost, že náhodně hrající hráč $\mathrm{O}$ zvolí po tahu do rohu 
právě tuto jedinou správnou odpověď, je:
\[
    \frac{1}{8}=0{,}125=12{,}5\,\%.
\]
Naopak pravděpodobnost, že zvolí některou ze zbývajících možností, 
po kterých může hráč $\mathrm{X}$ vynutit vítězství, je:
\[
    \frac{7}{8}=0{,}875=87{,}5\,\%.
\]
Tento výsledek však neznamená, 
že hráč $\mathrm{X}$ po chybném druhém tahu soupeře automaticky vyhraje. 
Znamená pouze, že od tohoto okamžiku existuje strategie, 
pomocí které může hráč $\mathrm{X}$ zvítězit bez ohledu na další tahy hráče $\mathrm{O}$.

Podívejme se nyní, proč je například po prvním tahu hráče $\mathrm{X}$ do rohu nutné, 
aby hráč $\mathrm{O}$ zahrál do středu hrací plochy.

Bez újmy na obecnosti předpokládejme, že hráč $\mathrm{X}$ začne tahem:
\[
    \mathrm{X}_1=p(1;1).
\]
Hráč $\mathrm{O}$ nyní udělá chybu a místo středu hrací plochy zahraje například:
\[
    \mathrm{O}_1=p(1;2).
\]
Hráč $\mathrm{X}$ následně zahraje do středu hrací plochy:
\[
    \mathrm{X}_2=p(2;2).
\]
V tomto okamžiku hrozí hráči $\mathrm{X}$ dokončení diagonály pomocí pole $p(3;3)$. 
Hráč $\mathrm{O}$ tedy musí tuto hrozbu zablokovat:
\[
    \mathrm{O}_2=p(3;3).
\]
Hráč $\mathrm{X}$ nyní zahraje:
\[
    \mathrm{X}_3=p(2;1).
\]
Dostáváme tak následující pozici:

\begin{figure}[h]
    \begin{center}
        \begin{tikzpicture}[scale=1.5]
            \draw[thick] (1,0) -- (1,3);
            \draw[thick] (2,0) -- (2,3);
            \draw[thick] (0,1) -- (3,1);
            \draw[thick] (0,2) -- (3,2);

            \node at (0.5,2.5) {\Large $\mathrm{X}$};
            \node at (1.5,2.5) {\Large $\mathrm{O}$};

            \node at (0.5,1.5) {\Large $\mathrm{X}$};
            \node at (1.5,1.5) {\Large $\mathrm{X}$};

            \node at (2.5,0.5) {\Large $\mathrm{O}$};

            \draw[red,thick,dashed] (2.5,1.5) circle (0.3);
            \draw[red,thick,dashed] (0.5,0.5) circle (0.3);
        \end{tikzpicture}
    \end{center}

    \caption{Pozice, ve které hráč $\mathrm{X}$ vytvořil dvě vítězné hrozby}
    \label{fig:TicTacToeVidlicka}
\end{figure}

Hráč $\mathrm{X}$ nyní hrozí vítězstvím na dvou různých místech současně.

Pokud v následujícím tahu zahraje:
\[
    \mathrm{X}=p(2;3),
\]
vytvoří vítěznou trojici ve druhém řádku.

Zároveň může zahrát:
\[
    \mathrm{X}=p(3;1),
\]
a vytvořit vítěznou trojici v prvním sloupci.

Hráč $\mathrm{O}$ může svým následujícím tahem zablokovat pouze jedno z těchto dvou polí.
Druhá vítězná možnost tedy zůstane volná a hráč $\mathrm{X}$ v následujícím tahu zvítězí.

Takovou pozici budeme nazývat \textit{vidličkou}.

Vidlička vzniká v okamžiku, kdy jeden hráč vytvoří současně dvě různé možnosti 
vítězství v následujícím tahu. Soupeř může běžným tahem zablokovat pouze jednu z nich,
a nemůže tedy zabránit vítězství.

V našem případě tedy po tahu:
\[
    \mathrm{X}_3=p(2;1)
\]
vznikly dvě vítězné možnosti:
\[
    p(2;3)
    \qquad\text{a}\qquad
    p(3;1).
\]
Vidlička je proto jedním z nejdůležitějších prvků strategie TicTacToe. 
Nestačí pouze sledovat, zda může některý z hráčů zvítězit v následujícím tahu. 
Je nutné také sledovat, zda některý tah nevytvoří dvě vítězné možnosti současně.

Na základě těchto úvah můžeme sestavit jednoduché pořadí priorit, 
podle kterého může hráč vybírat svůj následující tah.

\begin{enumerate}
    \item Pokud může hráč v aktuálním tahu zvítězit, provede vítězný tah.
    \item Pokud může soupeř zvítězit v následujícím tahu, hráč jeho vítěznou trojici zablokuje.
    \item Pokud může hráč vytvořit vidličku, vytvoří ji.
    \item Pokud může soupeř v následujícím tahu vytvořit vidličku, hráč jejímu vytvoření zabrání.
    \item Pokud je střed hrací plochy volný, hráč jej obsadí.
    \item Pokud soupeř obsadil některý roh a protější roh je volný, hráč obsadí protější roh.
    \item Pokud je volný některý z rohů, hráč obsadí jeden z nich.
    \item Pokud již není možné využít žádnou z předchozích možností, hráč obsadí některé volné pole uprostřed strany.
\end{enumerate}

Pořadí jednotlivých pravidel je důležité. Pokud například může hráč okamžitě zvítězit, 
nemá význam vytvářet vidličku nebo obsazovat střed hrací plochy. 
Stejně tak je nutné zabránit bezprostřednímu vítězství soupeře dříve, 
než se budeme zabývat vlastní vidličkou.

Poněkud složitější je zabránění soupeřově vidličce. 
Pokud existuje pouze jedno pole, jehož obsazením může soupeř vidličku vytvořit, 
stačí toto pole obsadit.

Mohou však nastat také situace, kdy má soupeř více možností vytvoření vidličky. 
V takovém případě nemusí být možné všechny tyto možnosti přímo zablokovat jediným tahem. 
Hráč proto musí vytvořit vlastní bezprostřední hrozbu vítězství. 
Soupeř je následně nucen tuto hrozbu zablokovat a nemůže svůj tah využít k vytvoření vidličky.

Toto pořadí pravidel představuje praktický způsob optimálního hraní TicTacToe. 
Pokud hráč tato pravidla správně dodržuje, soupeř si proti němu nemůže vynutit vítězství.

Pokud se stejnou strategií řídí oba hráči, 
žádný z nich nedokáže získat rozhodující výhodu a hra skončí remízou. 
To odpovídá výsledku, který jsme získali pomocí algoritmu minimax:
\[
    v(P_0)=0.
\]

Doposud jsme se na TicTacToe dívali převážně z matematického hlediska.
Nyní se můžeme podívat na to, jak předchozí výsledky využít přímo při hraní.

Jak již víme, pokud oba hráči hrají optimálně, skončí hra remízou.
Cílem optimální strategie tedy není nutně vždy zvítězit, 
protože proti bezchybně hrajícímu soupeři to není možné.
Cílem je hrát tak, abychom nikdy neprohráli 
a zároveň dokázali využít chybu soupeře k vítězství.

Nejprve se podíváme na strategii hráče $\mathrm{X}$, který provádí první tah.

\subsection*{Hra za $\mathrm{X}$}

Pro první tah hráče $\mathrm{X}$ máme po zohlednění symetrií tři možnosti:
střed hrací plochy, roh a střed některé ze stran.
Rozebereme proto všechny tři možnosti zvlášť.

\subsubsection*{Začátek ve středu}

Nejprve předpokládejme, že hráč $\mathrm{X}$ začne ve středu hrací plochy:
\[
    \mathrm{X}_1=p(2;2).
\]

Hráč $\mathrm{O}$ má nyní dva podstatně odlišné typy odpovědi.
Může zahrát do některého z rohů, nebo doprostřed některé ze stran.

Nejprve předpokládejme, že hráč $\mathrm{O}$ zahraje doprostřed strany.
Bez újmy na obecnosti můžeme zvolit:
\[
    \mathrm{O}_1=p(1;2).
\]

Hráč $\mathrm{X}$ nyní zahraje do jednoho z rohů sousedících s tokenem hráče $\mathrm{O}$.
Například:
\[
    \mathrm{X}_2=p(1;1).
\]

Dostáváme tedy pozici:

\begin{figure}[h]
    \begin{center}
        \begin{tikzpicture}[scale=1.5]
            \draw[thick] (1,0) -- (1,3);
            \draw[thick] (2,0) -- (2,3);
            \draw[thick] (0,1) -- (3,1);
            \draw[thick] (0,2) -- (3,2);

            \node at (0.5,2.5) {\Large $\mathrm{X}$};
            \node at (1.5,2.5) {\Large $\mathrm{O}$};
            \node at (1.5,1.5) {\Large $\mathrm{X}$};
        \end{tikzpicture}
    \end{center}

    \caption{Pozice po tahu hráče $\mathrm{O}$ doprostřed strany}
    \label{fig:TicTacToeStrategieXStredStrana}
\end{figure}

Hráč $\mathrm{X}$ nyní hrozí vytvořením vítězné diagonály:
\[
    p(1;1),\qquad p(2;2),\qquad p(3;3).
\]
Pokud hráč $\mathrm{O}$ pole $p(3;3)$ nezablokuje, 
zahraje na něj hráč $\mathrm{X}$ v následujícím tahu a zvítězí.

Hráč $\mathrm{O}$ je tedy nucen zahrát:
\[
    \mathrm{O}_2=p(3;3).
\]
Hráč $\mathrm{X}$ následně zahraje:
\[
    \mathrm{X}_3=p(2;1).
\]
Tím vzniknou dvě vítězné hrozby současně.
Hráč $\mathrm{X}$ může dokončit druhý řádek pomocí pole:
\[
    p(2;3),
\]
nebo první sloupec pomocí pole:
\[
    p(3;1).
\]
Hráč $\mathrm{O}$ může zablokovat pouze jednu z těchto možností.
Druhá zůstane volná a hráč $\mathrm{X}$ v následujícím tahu zvítězí.

Pokud tedy hráč $\mathrm{X}$ začne ve středu
a hráč $\mathrm{O}$ odpoví tahem doprostřed některé ze stran,
 může hráč $\mathrm{X}$ vynutit vítězství.

Nyní předpokládejme, že hráč $\mathrm{O}$ po prvním tahu hráče $\mathrm{X}$ zahraje do rohu.
Bez újmy na obecnosti zvolme:
\[
    \mathrm{O}_1=p(1;1).
\]
Hráč $\mathrm{X}$ zahraje do protějšího rohu:
\[
    \mathrm{X}_2=p(3;3).
\]
Dostáváme pozici:
\begin{figure}[h]
    \begin{center}
        \begin{tikzpicture}[scale=1.5]
            \draw[thick] (1,0) -- (1,3);
            \draw[thick] (2,0) -- (2,3);
            \draw[thick] (0,1) -- (3,1);
            \draw[thick] (0,2) -- (3,2);

            \node at (0.5,2.5) {\Large $\mathrm{O}$};
            \node at (1.5,1.5) {\Large $\mathrm{X}$};
            \node at (2.5,0.5) {\Large $\mathrm{X}$};
        \end{tikzpicture}
    \end{center}

    \caption{Pozice po tahu hráče $\mathrm{O}$ do rohu}
    \label{fig:TicTacToeStrategieXStredRoh}
\end{figure}
V tomto případě může hráč $\mathrm{O}$ správnou hrou zabránit vítězství hráče $\mathrm{X}$.
Aby si zachoval možnost remízy, musí nyní obsadit jeden ze dvou zbývajících rohů:
\[
    p(1;3)
    \qquad\text{nebo}\qquad
    p(3;1).
\]
Obě možnosti jsou vzhledem k symetrii stejné.
Předpokládejme tedy:
\[
    \mathrm{O}_2=p(1;3).
\]
Hráč $\mathrm{O}$ tím zároveň vytvořil hrozbu dokončení prvního řádku, 
a hráč $\mathrm{X}$ proto zahraje:
\[
    \mathrm{X}_3=p(1;2).
\]
Hráč $\mathrm{X}$ nyní hrozí dokončením druhého sloupce pomocí pole $p(3;2)$.
Hráč $\mathrm{O}$ je tedy při správné hře nucen tuto hrozbu zablokovat:
\[
    \mathrm{O}_3=p(3;2).
\]
V této situaci již žádný z hráčů nemůže při správné hře vynutit vítězství 
a hra může být dovedena k remíze.

Pokud však hráč $\mathrm{O}$ po tahu:
\[
    \mathrm{X}_2=p(3;3)
\]
nezahraje do jednoho ze zbývajících rohů, 
hráč $\mathrm{X}$ může jeho chybu využít a při správném pokračování si vynutit vítězství.

Pro začátek hráče $\mathrm{X}$ ve středu tedy dostáváme jednoduchý výsledek.

Pokud hráč $\mathrm{O}$ odpoví tahem doprostřed strany, 
hráč $\mathrm{X}$ zahraje do sousedního rohu a může si vynutit vítězství.

Pokud hráč $\mathrm{O}$ odpoví tahem do rohu,
hráč $\mathrm{X}$ zahraje do protějšího rohu.
Při správné další hře hráče $\mathrm{O}$ skončí hra remízou, 
ale chyba hráče $\mathrm{O}$ může opět umožnit hráči $\mathrm{X}$ vynutit vítězství.


\subsubsection*{Začátek v rohu}
Nyní předpokládejme, že hráč $\mathrm{X}$ začne v jednom z rohů hrací plochy.
Bez újmy na obecnosti můžeme zvolit:
\[
    \mathrm{X}_1=p(1;1).
\]
Hráč $\mathrm{O}$ může nyní odpovědět do středu, 
na některou ze stran nebo do některého ze zbývajících rohů.

Nejprve předpokládejme, že hráč $\mathrm{O}$ zahraje správně do středu:
\[
    \mathrm{O}_1=p(2;2).
\]
Hráč $\mathrm{X}$ následně zahraje do protějšího rohu:
\[
    \mathrm{X}_2=p(3;3).
\]
Dostáváme tedy pozici:
\begin{figure}[h]
    \begin{center}
        \begin{tikzpicture}[scale=1.5]
            \draw[thick] (1,0) -- (1,3);
            \draw[thick] (2,0) -- (2,3);
            \draw[thick] (0,1) -- (3,1);
            \draw[thick] (0,2) -- (3,2);

            \node at (0.5,2.5) {\Large $\mathrm{X}$};
            \node at (1.5,1.5) {\Large $\mathrm{O}$};
            \node at (2.5,0.5) {\Large $\mathrm{X}$};
        \end{tikzpicture}
    \end{center}

    \caption{Pozice po správné odpovědi hráče $\mathrm{O}$ do středu}
    \label{fig:TicTacToeStrategieXRohStred}
\end{figure}

Hráč $\mathrm{O}$ má nyní dvě podstatně odlišné možnosti.
Může zahrát doprostřed některé ze stran, nebo do jednoho ze dvou zbývajících rohů.

Pokud hráč $\mathrm{O}$ zahraje do jednoho ze zbývajících rohů, například:
\[
    \mathrm{O}_2=p(1;3),
\]
hráč $\mathrm{X}$ zahraje do posledního volného rohu:
\[
    \mathrm{X}_3=p(3;1).
\]
Tím vytvoří dvě vítězné hrozby současně.
Může dokončit třetí řádek pomocí pole $p(3;2)$ nebo první sloupec pomocí pole $p(2;1)$.

Hráč $\mathrm{O}$ může zablokovat pouze jednu z těchto možností 
a hráč $\mathrm{X}$ tedy v následujícím tahu zvítězí.

Pokud tedy hráč $\mathrm{O}$ po tazích:
\[
    \mathrm{X}_1=p(1;1),\qquad
    \mathrm{O}_1=p(2;2),\qquad
    \mathrm{X}_2=p(3;3)
\]
zahraje do některého ze zbývajících rohů, hráč $\mathrm{X}$ obsadí druhý volný roh 
a vytvoří vidličku.

Aby hráč $\mathrm{O}$ zabránil vítězství,
musí tedy ve svém druhém tahu zahrát doprostřed některé ze stran.

Předpokládejme například:
\[
    \mathrm{O}_2=p(1;2).
\]
Hráč $\mathrm{X}$ zahraje na protější stranu:
\[
    \mathrm{X}_3=p(3;2).
\]
Tím vytvoří hrozbu dokončení třetího řádku pomocí pole $p(3;1)$.
Hráč $\mathrm{O}$ je tedy nucen zahrát:
\[
    \mathrm{O}_3=p(3;1).
\]
Hráč $\mathrm{X}$ následně zahraje:
\[
    \mathrm{X}_4=p(1;3).
\]
Hráč $\mathrm{O}$ musí zablokovat třetí sloupec:
\[
    \mathrm{O}_4=p(2;3).
\]
Poslední volné pole je:
\[
    p(2;1),
\]
které obsadí hráč $\mathrm{X}$ a hra skončí remízou.

Dostáváme tedy úplnou posloupnost:
\[
\begin{aligned}
    &\mathrm{X}:p(1;1)
    \rightarrow
    \mathrm{O}:p(2;2)
    \rightarrow
    \mathrm{X}:p(3;3)
    \rightarrow\\
    &\rightarrow
    \mathrm{O}:p(1;2)
    \rightarrow
    \mathrm{X}:p(3;2)
    \rightarrow
    \mathrm{O}:p(3;1)
    \rightarrow\\
    &\rightarrow
    \mathrm{X}:p(1;3)
    \rightarrow
    \mathrm{O}:p(2;3)
    \rightarrow
    \mathrm{X}:p(2;1).
\end{aligned}
\]

Pokud hráč $\mathrm{O}$ všechny tyto tahy zahraje správně, skončí hra remízou.
Pokud však v některém okamžiku nezablokuje vznikající vítěznou hrozbu, 
může hráč $\mathrm{X}$ jeho chybu využít.

Nyní se vraťme k prvnímu tahu hráče $\mathrm{O}$ a předpokládejme, 
že místo středu zahraje doprostřed některé ze stran.

Například:
\[
    \mathrm{O}_1=p(1;2).
\]
Hráč $\mathrm{X}$ v takovém případě obsadí střed:
\[
    \mathrm{X}_2=p(2;2).
\]
Hráč $\mathrm{X}$ nyní hrozí dokončením diagonály pomocí pole $p(3;3)$.
Hráč $\mathrm{O}$ je proto nucen zahrát:
\[
    \mathrm{O}_2=p(3;3).
\]
Hráč $\mathrm{X}$ následně zahraje:
\[
    \mathrm{X}_3=p(2;1).
\]
Tím vytvoří vidličku, protože může zvítězit pomocí pole:
\[
    p(2;3)
\]
ve druhém řádku nebo pomocí pole:
\[
    p(3;1)
\]
v prvním sloupci.

Hráč $\mathrm{O}$ již nemůže obě možnosti současně zablokovat 
a hráč $\mathrm{X}$ si tedy vynutí vítězství.

Stejný princip platí i pro ostatní tahy hráče $\mathrm{O}$ doprostřed strany.
Hráč $\mathrm{X}$ nejprve obsadí střed hrací plochy a následně vytvoří dvě vítězné hrozby.

Zbývají případy, kdy hráč $\mathrm{O}$ odpoví na první tah hráče $\mathrm{X}$ tahem do rohu.

Pokud hráč $\mathrm{O}$ obsadí roh sousedící s prvním tokenem hráče $\mathrm{X}$, například:
\[
    \mathrm{O}_1=p(1;3),
\]
hráč $\mathrm{X}$ zahraje do rohu ležícího diagonálně proti tokenu hráče $\mathrm{O}$:
\[
    \mathrm{X}_2=p(3;1).
\]
Hráč $\mathrm{X}$ tím hrozí dokončením prvního sloupce pomocí pole $p(2;1)$.
Hráč $\mathrm{O}$ je tedy nucen tuto hrozbu zablokovat:
\[
    \mathrm{O}_2=p(2;1).
\]
Hráč $\mathrm{X}$ následně zahraje:
\[
    \mathrm{X}_3=p(3;3).
\]
Tím opět vytvoří dvě vítězné hrozby.
Může dokončit třetí řádek pomocí pole $p(3;2)$ nebo diagonálu pomocí pole $p(2;2)$.

Hráč $\mathrm{O}$ může zablokovat pouze jednu z těchto možností, 
a hráč $\mathrm{X}$ tedy zvítězí.

Pokud hráč $\mathrm{O}$ zahraje do rohu ležícího přímo 
proti prvnímu tokenu hráče $\mathrm{X}$:
\[
    \mathrm{O}_1=p(3;3),
\]
hráč $\mathrm{X}$ obsadí jeden ze zbývajících rohů, například:
\[
    \mathrm{X}_2=p(1;3).
\]
Hráč $\mathrm{X}$ nyní hrozí dokončením prvního řádku pomocí pole $p(1;2)$.
Hráč $\mathrm{O}$ je tedy nucen zahrát:
\[
    \mathrm{O}_2=p(1;2).
\]
Hráč $\mathrm{X}$ následně zahraje do posledního volného rohu:
\[
    \mathrm{X}_3=p(3;1).
\]
Tím vytvoří dvě vítězné možnosti:
\[
    p(2;1)
    \qquad\text{a}\qquad
    p(2;2).
\]
Hráč $\mathrm{O}$ může zablokovat pouze jednu z nich 
a hráč $\mathrm{X}$ si opět vynutí vítězství.

Pokud tedy hráč $\mathrm{X}$ začne v rohu, 
jedinou správnou první odpovědí hráče $\mathrm{O}$ je obsazení středu hrací plochy.

Pokud hráč $\mathrm{O}$ střed neobsadí, 
může hráč $\mathrm{X}$ při správném pokračování vždy vytvořit vidličku 
a vynutit si vítězství.

Pokud hráč $\mathrm{O}$ střed obsadí, zahraje hráč $\mathrm{X}$ do protějšího rohu.
Pokud následně hráč $\mathrm{O}$ zahraje do dalšího rohu, 
hráč $\mathrm{X}$ vytvoří vidličku obsazením posledního rohu.
Pokud hráč $\mathrm{O}$ správně zahraje doprostřed strany, 
může při bezchybné další hře dovést hru k remíze.

\subsubsection*{Začátek uprostřed strany}

Nakonec předpokládejme, že hráč $\mathrm{X}$ začne uprostřed některé ze stran hrací plochy.
Bez újmy na obecnosti zvolme horní stranu:
\[
    \mathrm{X}_1=p(1;2).
\]
Hráč $\mathrm{O}$ má nyní 8 možných odpovědí. 
Na rozdíl od předchozích případů však existuje několik různých tahů, 
kterými může zabránit vynucenému vítězství hráče $\mathrm{X}$.

První možností je obsazení středu hrací plochy:
\[
    \mathrm{O}_1=p(2;2).
\]
Hráč $\mathrm{X}$ následně zahraje do některého z horních rohů, například:
\[
    \mathrm{X}_2=p(1;1).
\]
Pokud oba hráči pokračují správně, nemůže si již hráč $\mathrm{X}$ vynutit vítězství, 
ale stále si může zajistit alespoň remízu.

Podobná situace nastane, pokud hráč $\mathrm{O}$ zahraje do 
jednoho z rohů sousedících s prvním tokenem hráče $\mathrm{X}$:
\[
    \mathrm{O}_1=p(1;1)
\]
nebo
\[
    \mathrm{O}_1=p(1;3).
\]
V takovém případě je pro hráče $\mathrm{X}$ výhodné obsadit střed hrací plochy:
\[
    \mathrm{X}_2=p(2;2).
\]
Ani v tomto případě si hráč $\mathrm{X}$ proti správně hrajícímu 
soupeři nemůže vynutit vítězství, ale může vždy dovést hru alespoň k remíze.

Poslední správnou odpovědí hráče $\mathrm{O}$ je tah doprostřed protější strany:
\[
    \mathrm{O}_1=p(3;2).
\]
Také v tomto případě může hráč $\mathrm{X}$ obsadit střed:
\[
    \mathrm{X}_2=p(2;2),
\]
a při správném pokračování si zachovat možnost remízy.

Hráč $\mathrm{O}$ tedy po prvním tahu:
\[
    \mathrm{X}_1=p(1;2)
\]
může zabránit vynucenému vítězství hráče $\mathrm{X}$ pomocí čtyř polí:
\[
    p(1;1),\qquad
    p(1;3),\qquad
    p(2;2),\qquad
    p(3;2).
\]
Zbývající čtyři tahy hráče $\mathrm{O}$ jsou chybné:
\[
    p(2;1),\qquad
    p(2;3),\qquad
    p(3;1),\qquad
    p(3;3).
\]
Pokud hráč $\mathrm{O}$ zahraje na některé z těchto polí, 
může si hráč $\mathrm{X}$ při správném pokračování vynutit vítězství.

Předpokládejme například, že hráč $\mathrm{O}$ zahraje:
\[
    \mathrm{O}_1=p(2;1).
\]
Hráč $\mathrm{X}$ v takovém případě zahraje do levého horního rohu:
\[
    \mathrm{X}_2=p(1;1).
\]
Hráč $\mathrm{X}$ tím vytvoří bezprostřední hrozbu dokončení prvního řádku pomocí pole:
\[
    p(1;3).
\]
Hráč $\mathrm{O}$ je tedy nucen tuto hrozbu zablokovat:
\[
    \mathrm{O}_2=p(1;3).
\]
Hráč $\mathrm{X}$ nyní obsadí střed hrací plochy:
\[
    \mathrm{X}_3=p(2;2).
\]
V tomto okamžiku vzniknou dvě vítězné hrozby současně. 
Hráč $\mathrm{X}$ může dokončit druhý sloupec pomocí pole:
\[
    p(3;2),
\]
nebo diagonálu pomocí pole:
\[
    p(3;3).
\]
Hráč $\mathrm{O}$ může zablokovat pouze jedno z těchto polí 
a hráč $\mathrm{X}$ tedy v následujícím tahu zvítězí.

Stejný princip funguje také při tahu:
\[
    \mathrm{O}_1=p(3;1).
\]
Hráč $\mathrm{X}$ opět zahraje:
\[
    \mathrm{X}_2=p(1;1).
\]
Hráč $\mathrm{O}$ musí zablokovat první řádek:
\[
    \mathrm{O}_2=p(1;3),
\]
a hráč $\mathrm{X}$ následně obsadí střed:
\[
    \mathrm{X}_3=p(2;2).
\]
Opět tak vytvoří vidličku na polích:
\[
    p(3;2)
    \qquad\text{a}\qquad
    p(3;3).
\]
Případy, kdy hráč $\mathrm{O}$ zahraje na pravou stranu hrací plochy, tedy na:
\[
    p(2;3)
    \qquad\text{nebo}\qquad
    p(3;3),
\]
jsou vzhledem k osové souměrnosti stejné. 
Hráč $\mathrm{X}$ v takovém případě zahraje do pravého horního rohu:
\[
    \mathrm{X}_2=p(1;3)
\]
a stejným postupem si vytvoří vidličku.

Pokud tedy hráč $\mathrm{X}$ začne uprostřed strany,
 může si vždy zajistit alespoň remízu.

Pokud hráč $\mathrm{O}$ odpoví do středu, 
do jednoho ze dvou sousedních rohů nebo doprostřed protější strany, 
může při správné další hře rovněž zajistit remízu.

Pokud však hráč $\mathrm{O}$ zahraje na některé ze zbývajících čtyř polí, 
může hráč $\mathrm{X}$ jeho chybu využít, vytvořit vidličku a vynutit si vítězství.

\subsection*{Hra za $\mathrm{O}$}

Nyní se podíváme na strategii hráče $\mathrm{O}$.
Ten je v jiné situaci než hráč $\mathrm{X}$, protože nezačíná 
a jeho první tah je vždy reakcí na první tah soupeře.

Hráč $\mathrm{O}$ si proti správně hrajícímu hráči $\mathrm{X}$ nemůže vynutit vítězství.
Může však vždy hrát tak, aby neprohrál a dovedl hru alespoň k remíze.
Pokud hráč $\mathrm{X}$ udělá chybu, 
může ji hráč $\mathrm{O}$ stejně jako v předchozích případech využít k vítězství.

\subsubsection*{Hráč $\mathrm{X}$ začne ve středu}

Předpokládejme, že hráč $\mathrm{X}$ začne ve středu hrací plochy:
\[
    \mathrm{X}_1=p(2;2).
\]
Hráč $\mathrm{O}$ by měl odpovědět tahem do některého z rohů.
Bez újmy na obecnosti zvolme:
\[
    \mathrm{O}_1=p(1;1).
\]
Další tah hráče $\mathrm{X}$ můžeme opět rozdělit podle toho, 
zda zahraje doprostřed strany nebo do některého ze zbývajících rohů.

Pokud hráč $\mathrm{X}$ zahraje doprostřed strany, například:
\[
    \mathrm{X}_2=p(1;2),
\]
musí hráč $\mathrm{O}$ zahrát na protější stranu:
\[
    \mathrm{O}_2=p(3;2).
\]

Obdobně pokud hráč $\mathrm{X}$ zahraje:
\[
    \mathrm{X}_2=p(2;1),
\]
odpoví hráč $\mathrm{O}$:
\[
    \mathrm{O}_2=p(2;3).
\]
Stejný princip platí i pro zbývající dvě strany.

Pokud hráč $\mathrm{X}$ zahraje do rohu sousedícího 
s prvním tokenem hráče $\mathrm{O}$, například:
\[
    \mathrm{X}_2=p(1;3),
\]
musí hráč $\mathrm{O}$ obsadit protější roh:
\[
    \mathrm{O}_2=p(3;1).
\]
Pokud naopak hráč $\mathrm{X}$ obsadí:
\[
    \mathrm{X}_2=p(3;1),
\]
zahraje hráč $\mathrm{O}$:
\[
    \mathrm{O}_2=p(1;3).
\]
Zvláštní případ nastane, pokud hráč $\mathrm{X}$ zahraje 
do rohu protějšího prvnímu tokenu hráče $\mathrm{O}$:
\[
    \mathrm{X}_2=p(3;3).
\]
V takovém případě musí hráč $\mathrm{O}$ obsadit jeden ze zbývajících dvou rohů:
\[
    \mathrm{O}_2=p(1;3)
\]
nebo
\[
    \mathrm{O}_2=p(3;1).
\]
Obě možnosti jsou vzhledem k symetrii stejné.

Pokud tedy hráč $\mathrm{X}$ začne ve středu, hráč $\mathrm{O}$ zahraje nejprve do rohu.
Při dalším tahu reaguje podle polohy druhého tokenu hráče $\mathrm{X}$ tak, 
aby zabránil vytvoření vidličky.
Při správném pokračování si tímto způsobem vždy zachová možnost remízy.

\subsubsection*{Hráč $\mathrm{X}$ začne v rohu}
Nyní předpokládejme, že hráč $\mathrm{X}$ začne v rohu hrací plochy.
Bez újmy na obecnosti zvolme:
\[
    \mathrm{X}_1=p(1;1).
\]
V tomto případě je správná první odpověď hráče $\mathrm{O}$ jednoznačná.
Musí obsadit střed hrací plochy:
\[
    \mathrm{O}_1=p(2;2).
\]
Pokud by hráč $\mathrm{O}$ zahrál kamkoliv jinam, 
může hráč $\mathrm{X}$ při správném pokračování vytvořit vidličku a vynutit si vítězství.

Nyní záleží na druhém tahu hráče $\mathrm{X}$.

Pokud hráč $\mathrm{X}$ zahraje doprostřed horní strany:
\[
    \mathrm{X}_2=p(1;2),
\]
hrozí dokončením prvního řádku.
Hráč $\mathrm{O}$ proto musí zahrát:
\[
    \mathrm{O}_2=p(1;3).
\]
Obdobně pokud hráč $\mathrm{X}$ zahraje:
\[
    \mathrm{X}_2=p(2;1),
\]
musí hráč $\mathrm{O}$ odpovědět:
\[
    \mathrm{O}_2=p(3;1).
\]
Pokud hráč $\mathrm{X}$ zahraje do sousedního rohu:
\[
    \mathrm{X}_2=p(1;3),
\]
musí hráč $\mathrm{O}$ zablokovat první řádek:
\[
    \mathrm{O}_2=p(1;2).
\]
Stejně tak při tahu:
\[
    \mathrm{X}_2=p(3;1)
\]
odpoví:
\[
    \mathrm{O}_2=p(2;1).
\]
Důležitý případ nastane, pokud hráč $\mathrm{X}$ obsadí protější roh:
\[
    \mathrm{X}_2=p(3;3).
\]
Hráč $\mathrm{O}$ v takovém případě nesmí zahrát do dalšího rohu.
Musí obsadit některé pole uprostřed strany, například:
\[
    \mathrm{O}_2=p(1;2).
\]
Stejně správné by bylo také:
\[
    p(2;1),\qquad p(2;3),\qquad p(3;2).
\]
Tím hráč $\mathrm{O}$ zabrání tomu, 
aby hráč $\mathrm{X}$ vytvořil vidličku obsazením některého ze zbývajících rohů.

Pokud hráč $\mathrm{X}$ ve svém druhém tahu zahraje doprostřed strany, 
která neleží vedle jeho prvního tokenu, například:
\[
    \mathrm{X}_2=p(2;3),
\]
může hráč $\mathrm{O}$ bezpečně zahrát například:
\[
    \mathrm{O}_2=p(1;3).
\]
Obdobně při tahu:
\[
    \mathrm{X}_2=p(3;2)
\]
může hráč $\mathrm{O}$ zahrát:
\[
    \mathrm{O}_2=p(3;1).
\]
Pokud tedy hráč $\mathrm{X}$ začne v rohu, 
musí hráč $\mathrm{O}$ jako první obsadit střed.
V dalších tazích pak blokuje vznikající vítězné možnosti 
a především zabrání vytvoření vidličky v případě, 
kdy hráč $\mathrm{X}$ obsadí dva protější rohy.

Při správném pokračování může hráč $\mathrm{O}$ vždy dovést hru alespoň k remíze.

\subsubsection*{Hráč $\mathrm{X}$ začne uprostřed strany}
Nakonec předpokládejme, že hráč $\mathrm{X}$ začne uprostřed některé ze stran.
Bez újmy na obecnosti zvolme:
\[
    \mathrm{X}_1=p(1;2).
\]
Hráč $\mathrm{O}$ má v tomto případě několik správných možností prvního tahu.
Pro praktickou hru je však nejjednodušší obsadit střed:
\[
    \mathrm{O}_1=p(2;2).
\]
Tím si hráč $\mathrm{O}$ zachová možnost remízy bez ohledu 
na další správné tahy hráče $\mathrm{X}$.

Pokud hráč $\mathrm{X}$ ve druhém tahu obsadí jeden z horních rohů, například:
\[
    \mathrm{X}_2=p(1;1),
\]
musí hráč $\mathrm{O}$ zablokovat první řádek:
\[
    \mathrm{O}_2=p(1;3).
\]
Obdobně při tahu:
\[
    \mathrm{X}_2=p(1;3)
\]
zahraje hráč $\mathrm{O}$:
\[
    \mathrm{O}_2=p(1;1).
\]
Pokud hráč $\mathrm{X}$ zahraje na levou stranu:
\[
    \mathrm{X}_2=p(2;1),
\]
může hráč $\mathrm{O}$ zahrát:
\[
    \mathrm{O}_2=p(1;1).
\]
Pokud hráč $\mathrm{X}$ zahraje na pravou stranu:
\[
    \mathrm{X}_2=p(2;3),
\]
může hráč $\mathrm{O}$ odpovědět:
\[
    \mathrm{O}_2=p(1;3).
\]
Pokud hráč $\mathrm{X}$ obsadí některý z dolních rohů, můžeme postupovat obdobně.
Například po tahu:
\[
    \mathrm{X}_2=p(3;1)
\]
může hráč $\mathrm{O}$ zahrát:
\[
    \mathrm{O}_2=p(1;1).
\]
Při tahu:
\[
    \mathrm{X}_2=p(3;3)
\]
může hráč $\mathrm{O}$ odpovědět:
\[
    \mathrm{O}_2=p(1;3).
\]
Zajímavý případ nastane, pokud hráč $\mathrm{X}$ obsadí střed protější strany:
\[
    \mathrm{X}_2=p(3;2).
\]
Tento tah je pro hráče $\mathrm{X}$ chybou.
Hráč $\mathrm{O}$ může nyní zahrát do některého z rohů, například:
\[
    \mathrm{O}_2=p(1;1),
\]
a při správném dalším pokračování si již může vynutit vítězství.

Pokud tedy hráč $\mathrm{X}$ začne uprostřed strany, 
je pro hráče $\mathrm{O}$ nejjednodušší odpovědí obsazení středu hrací plochy.
Při správném pokračování tím vždy zabrání vynucenému vítězství 
hráče $\mathrm{X}$ a zajistí si alespoň remízu.

Pokud navíc hráč $\mathrm{X}$ v některém z následujících tahů udělá chybu, 
může hráč $\mathrm{O}$ přejít z obranné strategie k vítězné.

\subsection*{Praktické shrnutí}
Předchozí strategie můžeme na závěr shrnout do několika základních situací.
Pokud hrajeme za $\mathrm{X}$, je z praktického hlediska výhodné začít v rohu.
Bez újmy na obecnosti tedy předpokládejme:
\[
    \mathrm{X}_1=p(1;1).
\]
Podle první odpovědi hráče $\mathrm{O}$ potom pokračujeme následovně.
\begin{figure}[h]
    \begin{center}
        \begin{tikzpicture}[scale=0.75]

            % O zahraje do středu
            \begin{scope}
                \draw[thick] (1,0) -- (1,3);
                \draw[thick] (2,0) -- (2,3);
                \draw[thick] (0,1) -- (3,1);
                \draw[thick] (0,2) -- (3,2);

                \node at (0.5,2.5) {\Large $\mathrm{X}_1$};
                \node at (1.5,1.5) {\Large $\mathrm{O}_1$};
                \node at (2.5,0.5) {\Large $\mathrm{X}_2$};

                \node[align=center] at (1.5,-0.75)
                    {\small $\mathrm{O}$ hraje střed\\ \small $\mathrm{X}$ hraje protější roh};
            \end{scope}

            % O zahraje na stranu
            \begin{scope}[xshift=4.5cm]
                \draw[thick] (1,0) -- (1,3);
                \draw[thick] (2,0) -- (2,3);
                \draw[thick] (0,1) -- (3,1);
                \draw[thick] (0,2) -- (3,2);

                \node at (0.5,2.5) {\Large $\mathrm{X}_1$};
                \node at (1.5,2.5) {\Large $\mathrm{O}_1$};
                \node at (1.5,1.5) {\Large $\mathrm{X}_2$};

                \node[align=center] at (1.5,-0.75)
                    {\small $\mathrm{O}$ hraje na stranu\\ \small $\mathrm{X}$ hraje střed};
            \end{scope}

            % O zahraje do sousedního rohu
            \begin{scope}[xshift=9cm]
                \draw[thick] (1,0) -- (1,3);
                \draw[thick] (2,0) -- (2,3);
                \draw[thick] (0,1) -- (3,1);
                \draw[thick] (0,2) -- (3,2);

                \node at (0.5,2.5) {\Large $\mathrm{X}_1$};
                \node at (2.5,2.5) {\Large $\mathrm{O}_1$};
                \node at (0.5,0.5) {\Large $\mathrm{X}_2$};

                \node[align=center] at (1.5,-0.75)
                    {\small $\mathrm{O}$ hraje vedlejší roh\\ \small $\mathrm{X}$ hraje protější roh};
            \end{scope}

            % O zahraje do protějšího rohu
            \begin{scope}[xshift=13.5cm]
                \draw[thick] (1,0) -- (1,3);
                \draw[thick] (2,0) -- (2,3);
                \draw[thick] (0,1) -- (3,1);
                \draw[thick] (0,2) -- (3,2);

                \node at (0.5,2.5) {\Large $\mathrm{X}_1$};
                \node at (2.5,0.5) {\Large $\mathrm{O}_1$};
                \node at (2.5,2.5) {\Large $\mathrm{X}_2$};

                \node[align=center] at (1.5,-0.75)
                    {\small $\mathrm{O}$ hraje protější roh\\ \small $\mathrm{X}$ hraje volný roh};
            \end{scope}

        \end{tikzpicture}
    \end{center}

    \caption{Základní odpovědi hráče $\mathrm{X}$ po zahájení v rohu}
    \label{fig:TicTacToeShrnutiX}
\end{figure}

Pokud hráč $\mathrm{O}$ obsadí střed, může při bezchybné další hře dovést hru k remíze.
Hráč $\mathrm{X}$ proto zahraje do protějšího rohu a čeká na případnou další chybu soupeře.

Pokud hráč $\mathrm{O}$ zahraje na některou ze stran, hráč $\mathrm{X}$ obsadí střed.
Při správném pokračování si potom může vynutit vítězství.

Pokud hráč $\mathrm{O}$ zahraje do některého z ostatních rohů, 
hráč $\mathrm{X}$ pokračuje v jiném rohu tak, aby si připravil vznik vidličky.
Také v tomto případě může při správném pokračování využít chybu 
hráče $\mathrm{O}$ k vynucenému vítězství.

Pro hráče $\mathrm{O}$ je praktická strategie ještě jednodušší.
Jeho první tah závisí pouze na tom, kde zahájil hráč $\mathrm{X}$.

\begin{figure}[h]
    \begin{center}
        \begin{tikzpicture}[scale=1]

            % X začne ve středu
            \begin{scope}
                \draw[thick] (1,0) -- (1,3);
                \draw[thick] (2,0) -- (2,3);
                \draw[thick] (0,1) -- (3,1);
                \draw[thick] (0,2) -- (3,2);

                \node at (1.5,1.5) {\Large $\mathrm{X}_1$};
                \node at (0.5,2.5) {\Large $\mathrm{O}_1$};

                \node[align=center] at (1.5,-0.5)
                    {$\mathrm{X}$ začne ve středu\\$\mathrm{O}$ hraje do rohu};
            \end{scope}

            % X začne v rohu
            \begin{scope}[xshift=5cm]
                \draw[thick] (1,0) -- (1,3);
                \draw[thick] (2,0) -- (2,3);
                \draw[thick] (0,1) -- (3,1);
                \draw[thick] (0,2) -- (3,2);

                \node at (0.5,2.5) {\Large $\mathrm{X}_1$};
                \node at (1.5,1.5) {\Large $\mathrm{O}_1$};

                \node[align=center] at (1.5,-0.5)
                    {$\mathrm{X}$ začne v rohu\\$\mathrm{O}$ hraje do středu};
            \end{scope}

            % X začne na straně
            \begin{scope}[xshift=10cm]
                \draw[thick] (1,0) -- (1,3);
                \draw[thick] (2,0) -- (2,3);
                \draw[thick] (0,1) -- (3,1);
                \draw[thick] (0,2) -- (3,2);

                \node at (1.5,2.5) {\Large $\mathrm{X}_1$};
                \node at (1.5,1.5) {\Large $\mathrm{O}_1$};

                \node[align=center] at (1.5,-0.5)
                    {$\mathrm{X}$ začne na straně\\$\mathrm{O}$ hraje do středu};
            \end{scope}

        \end{tikzpicture}
    \end{center}

    \caption{Doporučené první odpovědi hráče $\mathrm{O}$}
    \label{fig:TicTacToeShrnutiO}
\end{figure}

Pokud hráč $\mathrm{X}$ začne ve středu, zahraje hráč $\mathrm{O}$ do některého z rohů.

Pokud hráč $\mathrm{X}$ začne v rohu, zahraje hráč $\mathrm{O}$ do středu.

Pokud hráč $\mathrm{X}$ začne uprostřed strany, 
je pro hráče $\mathrm{O}$ opět nejjednodušší bezpečnou odpovědí obsazení středu.

Při dodržování těchto zahajovacích principů 
a správné reakci na další hrozby si může každý 
z hráčů proti optimálně hrajícímu soupeři zajistit alespoň remízu.
Hráč, který jako první udělá chybu umožňující vznik vidličky 
nebo nezablokuje vítěznou hrozbu, 
však může soupeři umožnit přejít z remízy k vynucenému vítězství.

\section*{Závěr}

Na první pohled je TicTacToe velmi jednoduchá hra s pouze devíti poli a několika málo pravidly.
Při podrobnějším pohledu však zjistíme, že i takto jednoduchá hra nabízí poměrně zajímavý kombinatorický problém.

Nejprve jsme zjistili počet všech možných posloupností tahů a následně jsme pomocí symetrií hrací plochy odstranili hry, které se od sebe liší pouze otočením nebo osovou souměrností.
Po zohlednění toho, že hra může skončit již od pátého tahu, jsme získali celkem:
\[
    31\,896
\]
různých průběhů hry, pokud hry lišící se pouze symetrií považujeme za stejné.

Dále jsme se nezabývali pouze pořadím jednotlivých tahů, ale také pozicemi, které mohou během hry vzniknout.
Bez zohlednění symetrií může během hry nastat:
\[
    5\,478
\]
legálních pozic včetně prázdné hrací plochy.

Po zohlednění otočení a osových souměrností se tento počet sníží na:
\[
    765
\]
navzájem různých pozic.

Pomocí herního stromu a algoritmu minimax jsme dále zjistili, že při optimální hře obou hráčů nemůže ani jeden z nich vynutit vítězství.
Výsledkem optimální hry je tedy vždy:
\[
    \mathrm{X}\;-\;\mathrm{O}=\text{remíza}.
\]

To však neznamená, že jsou všechny tahy stejně dobré.
Pokud některý z hráčů udělá chybu, může soupeř v některých případech vytvořit vidličku a přejít z pozice vedoucí k remíze k vynucenému vítězství.

Na závěr jsme proto sestavili praktickou strategii pro oba hráče.
Pokud se jí hráč správně řídí, nemůže proti němu soupeř vynutit vítězství.
V nejhorším případě tedy hra skončí remízou, zatímco chybu soupeře lze využít k vítězství.

TicTacToe je tak příkladem hry, která je pravidly velmi jednoduchá, ale přesto na ní lze ukázat kombinatoriku, symetrie, herní stromy, algoritmus minimax i základní principy optimální strategie.

\end{document}